\documentclass[11pt,a4paper]{article}

% ==================================================
% GÓI THƯ VIỆN CHUẨN
% ==================================================
\usepackage[margin=2cm]{geometry}
\usepackage{amsmath,amssymb,amsthm,amsfonts,dsfont}
\usepackage[T1]{fontenc}
\usepackage[utf8]{inputenc}
\usepackage{enumitem}
\usepackage[most]{tcolorbox}
\usepackage{dsfont}
\usepackage[OT1]{fontenc}

% Cấu hình khoảng cách đoạn
\setlength{\parindent}{0pt}
\setlength{\parskip}{6pt}
\setlist{noitemsep, topsep=2pt}

% ==================================================
% ĐỊNH NGHĨA KHUNG ĐEN TRẮNG TỐI GIẢN (COMPONENTS)
% ==================================================

% 1. Khung Định lý (Theorem Box)
\newtcolorbox{thmbox}[1][]{
  colback=white,
  colframe=black,
  fonttitle=\bfseries,
  coltitle=black,
  attach title to upper={\ --\ },
  title={Theorem #1},
  arc=0mm,
  boxrule=0.8pt
}

% 2. Khung Định nghĩa (Definition Box)
\newtcolorbox{defbox}[1][]{
  colback=white,
  colframe=black,
  fonttitle=\bfseries,
  coltitle=black,
  attach title to upper={\ --\ },
  title={Definition #1},
  arc=0mm,
  boxrule=0.8pt
}

% 3. Khung Chú ý / Lưu ý (Note Box)
\newtcolorbox{notebox}[1][]{
  colback=white,
  colframe=black,
  fonttitle=\bfseries,
  coltitle=black,
  attach title to upper={\ --\ },
  title={Note #1},
  arc=0mm,
  boxrule=0.5pt,
  sharp corners
}

% 4. Khung Ví dụ (Example Box)
\newtcolorbox{exbox}[1][]{
  colback=white,
  colframe=black!60,
  fonttitle=\bfseries,
  coltitle=black,
  attach title to upper={\ --\ },
  title={Example #1},
  arc=0mm,
  boxrule=0.5pt
}

% ==================================================
% BẮT ĐẦU TÀI LIỆU
% ==================================================
\begin{document}

% --- TIÊU ĐỀ HANDBOOK ---
\begin{center}
    {\Huge \bfseries Convex Optimization Handbook} \\[0.2cm]
    {\large My cheatsheet for convex optimization} \\
    \vspace{0.2cm}
\end{center}

\vspace{0.2cm}

% ==================================================
% NỘI DUNG CHÍNH
% ==================================================

\section{Convexity}

\begin{thmbox}[(Weierstrass)]
Given $\Omega \subseteq \mathds{E}$ non-empty and compact, $f: \Omega \rightarrow \mathds{R}$ is continuous. Then:
\[
\exists x^* \in \Omega : f(x^*) = \min_{x \in \Omega} f(x)
\]

\textit{Interpretation}: Continuous functions always have minimizer on a non-empty and compact set. 
\end{thmbox}

\begin{defbox}[(Convex set)]
Given $C \in \mathds{E}$. $C$ is convex if:
\[
\forall x,y \in C, \theta \in [0,1] \Rightarrow \theta x+ (1-\theta y) \in C
\]
\textit{Interpretation}: The line connect any 2 points in the set must lie completely in the set.
\end{defbox}

\begin{defbox}[(Convex function)]
$\forall x,y \in C$, $C$ is convex, $\forall \theta \in [0,1], f: C \rightarrow \mathds{R}$. $f$ is convex if:
\[
f(\theta x + (1- \theta) y) \le \theta f(x) + (1 - \theta) f(y)
\]    
\end{defbox}

\begin{notebox}
With the strictly convex case, the inequation above won't have the equal sign.

\end{notebox}

\begin{notebox}
    Convex functions have to be defined on a convex set.
\end{notebox}

\begin{thmbox}[(Local minima are global for convex)]
$\emptyset \neq \Omega \subseteq \mathds{E}$ is convex, $f: \Omega \rightarrow \mathds{R}$ is convex. With $x^* \in \Omega$. If $x^*$ is a local minimum, then $x^*$ is a global minimum.
\end{thmbox}

\begin{thmbox}[(The non-existence of feasible descent direction of global minima in a convex set)]
$\emptyset \neq \Omega \subseteq \mathds{E}$ is convex, $f: \mathds{E} \rightarrow \mathds{R}$ is differentiable. $x^*$ is global minima of $f$ over $\Omega$. Then:
\[
\forall x \in \Omega: \langle\nabla f(x^*), x - x^*\rangle \ge 0 
\]

\textit{Interpretation}: The theorem just want to tell us that if you are at the global minimum of a convex set, there won't be any feasible descent directions.
\end{thmbox}

\begin{thmbox}[(The "tangent")]
$\emptyset \neq \Omega \subseteq \mathds{E}$ is convex, $f: \mathds{E} \rightarrow \mathds{R}$ is differentiable and convex on $\Omega$. Then, $\forall x,y \in \Omega$:
\[
f(y) \ge f(x) + \langle\nabla f(x), y-x\rangle
\]
\end{thmbox}

\begin{thmbox}[(First-order related property for convexity)]
$\emptyset \neq \Omega \subseteq \mathds{E}$ is convex, $f: \mathds{E} \rightarrow \mathds{R}$ is differentiable and convex on $\Omega$. With $x^* \in \Omega$. The following is equivalent:
\begin{itemize}
    \item \textbf{1.}\[f(x^*) \le f(x) \quad \forall x \in \Omega\]
    \item \textbf{2.}\[\langle \nabla f(x^*), x-x^*\rangle \quad \forall x \in \Omega\]
\end{itemize}
\end{thmbox}

\section{The geometry behind the "globalness" of the locals in convex functions}

\begin{defbox}[(Extended real line)]
$\mathds{R} \cup \{+\infty\}$ is the one-sided extended real line, with $a \le +\infty$, we have the convention:
\[a + (+\infty) = +\infty\]
\[\lambda(+\infty) = +\infty \quad \forall \lambda > 0\]
\[0(+\infty) := 0\]
\[\inf \emptyset = +\infty\]
\end{defbox}

\begin{defbox}[(Effective domain)]
$f: \mathds{E} \rightarrow \mathds{R} \cup \{+\infty\}$, the effective domain is:
\[
domf = \{x \in \mathds{E}: f(x) < +\infty\}
\]
\end{defbox}

\begin{defbox}[(Convex extended real-valued function)]
$f: \mathds{E} \rightarrow \mathds{R} \cup \{+\infty\}$ is convex if:
\[
\forall x,y \in \mathds{E}, \theta \in [0,1] \quad f(\theta x + (1 - \theta)y) \le \theta f(x) + (1-\theta)f(y)
\]
\end{defbox}

\begin{notebox}
The definition for convex function above is the same as Part 1. We just changed the input's domain in to $\mathds{E}$ instead of the feasible set $C$. You can transform the Part 1 definition into the above definition by define $\Tilde{f}(x) = f(x) \quad \forall x \in C$, else $+\infty$. Then $f$ is convex iff $\Tilde{f}$ is convex.
\end{notebox}

\begin{defbox}[(Epigraph)]
For $f: \mathds{E} \rightarrow \mathds{R} \cup \{+\infty\}$, the epigraph is:
\[
epi(f) = \{(x,t) \in \mathds{E} \times \mathds{R}: t \le f(x)\}
\]
\end{defbox}

\begin{thmbox}[(Epigraph links with convexity of a function)]
$f$ convex $\Leftrightarrow$ $epi(f)$ is convex
\end{thmbox}


\begin{defbox}[(Closed halfspace)]
For $\xi \in E^* \setminus \{0\}$ and $b \in \mathbb{R}$:
\[
H(\xi, b) := \{\, x \in E : \langle \xi, x \rangle \le b \,\}.
\]
\end{defbox}

\begin{defbox}[(Interior)]
\[
\operatorname{int}(S) := \{\, x \in S : \exists\, r > 0,\ B(x, r) \subseteq S \,\}
\]
\textit{Interpretation}: Everything of $S$, except its borders.
\end{defbox}

\begin{defbox}[(Closure)]
 \[
 \operatorname{cl}(S) := \bigcap \{\, F : F \text{ closed},\ S \subseteq F \,\}
 \]
 \textit{Interpretation}: Everything of $S$, plus with missing borders.
\end{defbox}

\newpage
% ==================================================
% KHO COMPONENT (CHỈ DÙNG ĐỂ COPY LÊN TRÊN)
% ==================================================
\section*{comp}
\hrule
\vspace{0.5cm}

% --- COMPONENT 1: THM BOX ---
% \begin{thmbox}[Tên Định Lý]
% Nội dung định lý...
% \end{thmbox}

% --- COMPONENT 2: DEF BOX ---
% \begin{defbox}[Tên Định Nghĩa]
% Nội dung định nghĩa...
% \end{defbox}

% --- COMPONENT 3: NOTE BOX ---
% \begin{notebox}[Tên Chú Ý]
% Nội dung chú ý...
% \end{notebox}

% --- COMPONENT 4: EXAMPLE BOX ---
% \begin{exbox}[Tên Ví Dụ]
% Nội dung ví dụ...
% \end{exbox}

% --- COMPONENT 5: HEADING / SECTION ---
% \section{Tên Chương/Mục Chính}
% \subsection{Tên Mục Phụ}

% --- COMPONENT 6: LIST (DANH SÁCH) ---
% \begin{itemize}
%     \item \textbf{Ý thứ nhất:} Chi tiết...
%     \item \textbf{Ý thứ hai:} Chi tiết...
% \end{itemize}

\end{document}